\newcounter{english}
\documentclass{article}

% packages
\usepackage{amsmath, amsthm, thmtools, amsfonts, amssymb, luacode, catchfile, tikzducks, hyperref, ifthen}
\ifcsname c@kobocompile\endcsname
	\usepackage[a5paper, total={1072pt, 1448pt}, margin=10pt, includeheadfoot]{geometry} % set page margins
\else
	\usepackage[a4paper, margin=50pt, includeheadfoot]{geometry}
\fi
\usepackage[shortlabels]{enumitem}
\usepackage[skip=3pt, indent=0pt]{parskip}

% language
\usepackage[bidi=basic, layout=tabular, provide=*]{babel}
\ifcsname c@english\endcsname
	\babelprovide[main, import]{english}
\else
	\babelprovide[main, import]{hebrew}
	\babelprovide{rl}
\fi
%\babelfont{rm}{Libertinus Serif}
\babelfont{rm}[Renderer=Harfbuzz]{Libertinus Serif}
\babelfont{sf}{Libertinus Sans}
\babelfont{tt}{Libertinus Mono}

% style
\AddToHook{cmd/section/before}{\clearpage}	% Add line break before section
\linespread{1.3}
\setcounter{secnumdepth}{0}		% Remove default number tags from sections, this won't do well with theorems
\AtBeginDocument{\setlength{\belowdisplayskip}{3pt}}
\AtBeginDocument{\setlength{\abovedisplayskip}{3pt}}
\graphicspath{ {../images/} }

% operators
\DeclareMathOperator\cis{cis}
\DeclareMathOperator\Sp{Sp}
\DeclareMathOperator\tr{tr}
\DeclareMathOperator\im{Im}
\DeclareMathOperator\re{Re}
\DeclareMathOperator\diag{diag}
\DeclareMathOperator*\lowlim{\underline{lim}}
\DeclareMathOperator*\uplim{\overline{lim}}
\DeclareMathOperator\rng{rng}
\DeclareMathOperator\Sym{Sym}
\DeclareMathOperator\Arg{Arg}
\DeclareMathOperator\Log{Log}
\DeclareMathOperator\dom{dom}
\DeclareMathOperator\supp{Supp}
\DeclareMathOperator\var{Var}
\DeclareMathOperator\cov{Cov}

% commands
%\renewcommand\qedsymbol{\textbf{מש''ל}}
%\renewcommand\qedsymbol{\fbox{\emoji{lizard}}}
\newcommand{\Aa}[0]{\mathcal{A}}
\newcommand{\Bb}[0]{\mathcal{B}}
\newcommand{\CC}[0]{\mathbb{C}}
\newcommand{\Cc}[0]{\mathcal{C}}
\newcommand{\EE}[0]{\mathbb{E}}
\newcommand{\FF}[0]{\mathbb{F}}
\newcommand{\Ff}[0]{\mathcal{F}}
\newcommand{\Ii}[0]{\mathcal{I}}
\newcommand{\Gg}[0]{\mathcal{G}}
\newcommand{\Ll}[0]{\mathcal{L}}
\newcommand{\Mm}[0]{\mathcal{M}}
\newcommand{\NN}[0]{\mathbb{N}}
\newcommand{\Nn}[0]{\mathcal{N}}
\newcommand{\PP}[0]{\mathbb{P}}
\newcommand{\Pp}[0]{\mathcal{P}}
\newcommand{\QQ}[0]{\mathbb{Q}}
\newcommand{\RR}[0]{\mathbb{R}}
\newcommand{\Rr}[0]{\mathcal{R}}
\newcommand{\Ss}[0]{\mathcal{S}}
\newcommand{\TT}[0]{\mathbb{T}}
\newcommand{\Uu}[0]{\mathcal{U}}
\newcommand{\Vv}[0]{\mathcal{V}}
\newcommand{\Ww}[0]{\mathcal{W}}
\newcommand{\ZZ}[0]{\mathbb{Z}}
\newcommand{\acts}[0]{\circlearrowright}
\newcommand{\explain}[2] {
	\begin{flalign*}
		 && \text{#2} && \text{#1}
	\end{flalign*}
}
\newcommand{\maketitleprint}[0]{ \begin{center}
	%\begin{tikzpicture}[scale=3]
	%	\duck[graduate=gray!20!black, tassel=red!70!black]
	%\end{tikzpicture}	
	\includegraphics[width=6cm]{cover}
\end{center}
}

% theorem commands
\newtheoremstyle{c_remark}
	{}	% Space above
	{}	% Space below
	{}% Body font
	{}	% Indent amount
	{\bfseries}	% Theorem head font
	{}	% Punctuation after theorem head
	{.5em}	% Space after theorem head
	{\thmname{#1}\thmnumber{ #2}\thmnote{ \normalfont{\text{(#3)}}}}	% head content
\newtheoremstyle{c_definition}
	{3pt}	% Space above
	{3pt}	% Space below
	{}% Body font
	{}	% Indent amount
	{\bfseries}	% Theorem head font
	{}	% Punctuation after theorem head
	{.5em}	% Space after theorem head
	{\thmname{#1}\thmnumber{ #2}\thmnote{ \normalfont{\text{(#3)}}}}	% head content
\newtheoremstyle{c_plain}
	{3pt}	% Space above
	{3pt}	% Space below
	{\itshape}% Body font
	{}	% Indent amount
	{\bfseries}	% Theorem head font
	{}	% Punctuation after theorem head
	{.5em}	% Space after theorem head
	{\thmname{#1}\thmnumber{ #2}\thmnote{ \text{(#3)}}}	% head content

\ifcsname c@english\endcsname
	\theoremstyle{plain}
	\newtheorem{theorem}{Theorem}[section]
	\newtheorem{lemma}[theorem]{Lemma}
	\newtheorem{proposition}[theorem]{Proposition}
	\newtheorem*{proposition*}{Proposition}
	%\newtheorem{corollary}[theorem]{אין חלופה עברית}

	\theoremstyle{definition}
	\newtheorem{definition}[theorem]{Definition}
	\newtheorem*{definition*}{Definition}
	\newtheorem{example}{Example}[section]
	\newtheorem{exercise}{Exercise}[section]

	\theoremstyle{remark}
	\newtheorem*{remark}{Remark}
	\newtheorem*{solution}{Solution}
	\newtheorem{conclusion}[theorem]{Conclusion}
	\newtheorem{notation}[theorem]{Notation}
\else
	\theoremstyle{c_plain}
	\newtheorem{theorem}{משפט}[section]
	\newtheorem{lemma}[theorem]{למה}
	\newtheorem{proposition}[theorem]{טענה}
	\newtheorem*{proposition*}{טענה}
	%\newtheorem{corollary}[theorem]{אין חלופה עברית}

	\theoremstyle{c_definition}
	\newtheorem{definition}[theorem]{הגדרה}
	\newtheorem*{definition*}{הגדרה}
	\newtheorem{example}{דוגמה}[section]
	\newtheorem{exercise}{תרגיל}[section]

	\theoremstyle{c_remark}
	\newtheorem*{remark}{הערה}
	\newtheorem*{solution}{פתרון}
	\newtheorem{conclusion}[theorem]{מסקנה}
	\newtheorem{notation}[theorem]{סימון}
\fi

% Questions related commands
\newcounter{question}
\setcounter{question}{1}
\newcounter{sub_question}
\setcounter{sub_question}{1}

\ifcsname c@english\endcsname
	\newcommand{\question}[1][0]{
		\ifthenelse{#1 = 0}{}{\setcounter{question}{#1}}
		\section{Question \arabic{question}}
		\addtocounter{question}{1}
		\setcounter{sub_question}{1}
	}

	\newcommand{\subquestion}[1][0]{
		\ifthenelse{#1 = 0}{}{\setcounter{sub_question}{#1}}
		\subsection{Part \alph{sub_question}}
		\addtocounter{sub_question}{1}
	}
\else
	\newcommand{\question}[1][0]{
		\ifthenelse{#1 = 0}{}{\setcounter{question}{#1}}
		\section{שאלה \arabic{question}}
		\addtocounter{question}{1}
		\setcounter{sub_question}{1}
	}

	\newcommand{\subquestion}[1][0]{
		\ifthenelse{#1 = 0}{}{\setcounter{sub_question}{#1}}
		\subsection{סעיף \localecounter{letters.gershayim}{sub_question}}
		\addtocounter{sub_question}{1}
	}
\fi

% import lua and start of document
\directlua{common = require ('../common')}

\GetEnv{AUTHOR}

% headers
\author{\AUTHOR}
\date\today

\title{Exercise 1 Answer Sheet --- Logic Theory (2), 80424}

\begin{document}
\maketitle
\maketitleprint{}

\question{}
Let us assume that $D$ is an ultrafilter on given set $X$.

\subquestion{}
Suppose that there is some $s \in D$ such that $|s| < \omega$, we will show that $D$ is principal.
\begin{proof}
	We will assume that $s$ is minimal in relation to cardinality in $D$, otherwise let us choose the minimal such $s$ by using the well ordering principle.
	There is a unique such $s \in D$, otherwise $|s| = |s'|$ and $|s \cap s'| < |s|$, a contradiction to the minimality of the cardinality of $s$. \\
	For every $x \in D$, there is some $y \in D$ such that $y \subseteq x \cap s$, but by $s$ minimality, $y = x \cap s$, hence $s \subseteq x$.
	In the other direction, if $s \subseteq x$, then $x \in D$ by filter properties. \\
	Therefore $D = \{ x \in \Pp(X) \mid s \subseteq x \}$, $D$ is a principle ultrafilter defined by $s$.
\end{proof}

\subquestion{}
We will show that if $D \subseteq \Pp(I)$ is a principal ultrafilter, $\langle \Mm_i \mid i \in I \rangle$ is a sequence of $L$-structures, then $\Nn = \prod_{i \in I} \Mm_i / D \cong M_i$ for some $i \in I$.
\begin{proof}
	Let $s \subseteq I$ be the defining set for $D$, and let us fix $i \in s$.
	Let $f : N \to M_i$ such that $f({[g]}_D) = g(i)$.
	$f$ is well defined, as every for every $h, h' \in [h]$, it is follows that $h \restriction s = h' \restriction s$ by $D$'s definition as principle.

	We will show that $f$ is an isomorphism between $\Nn$ and $\Mm_i$.
	$L$-constants are being preserved directly from definition of ultra-products.
	Assume $F \in \operatorname{Func}_{L, n}$, and $t_0, \dots, t_{n - 1} \in \operatorname{term}_L$, then,
	\[
		f(F^\Nn(t_0^\Nn, \dots, t_{n - 1}^\Nn))
		= f([\{ F^{\Mm_j}(t_0^{\Mm_j}, \dots, t_{n - 1}^{\Mm_j}) \mid j \in I \}])
		= F^{\Mm_i}(t_0^{\Mm_i}, \dots, t_{n - 1}^{\Mm_i})
	\]
	meaning $f$ preserves terms.
	Let $R \in \operatorname{rel}_{L, n}$, and let $t_0, \dots, t_{n - 1} \in \operatorname{term}_L$, then,
	\[
		\Nn \models R(t_0, \dots, t_{n - 1})
		\iff \{ j \in I \mid \langle t_0^{\Mm_j}, \dots, t_{n - 1}^{\Mm_j} \rangle \in R^{\Mm_j} \} \in D
		\implies \langle t_0^{\Mm_j}, \dots, t_{n - 1}^{\Mm_j} \rangle \in R^{\Mm_i}
	\]
	We conclude that $f$ is indeed a structures isomorphism.
\end{proof}

\question{}
We will show that if $F$ is a filter on a set $X$ such that if $s \in F$ then $s$ is infinite, then there is a non-principle ultrafilter extending $F$.
\begin{proof}
	Initially, we claim that there is no $x \in F$ such that $|x|, |x^C| \ge \omega$, when $x^C = X \setminus x$.
	If there is one, then $\emptyset \subseteq x \cap x^C = \emptyset$ is in $F$, in contradiction to $F$ being a filter.
	We define the ultrafilter $U = \{ x \subseteq X \mid |x| \ge \omega \}$, this is indeed an ultrafilter as noted in lecture.
	it is sufficient to show that $F \subseteq U$.
	Indeed, if $s \in F$ then $|s| \ge \omega$, implies that $s \in U$.

	We found an ultrafilter extending $F$, it is required to prove that $U$ is not a principle ultrafilter as well.
	We will assume otherwise, then it is follows that there is $S \subseteq X$ such that $\forall x, S \subseteq x \iff x \in U$.
	Let $a \in S$ be an element, then $b = X \setminus \{a\}$ is a set such that $|b| \ge \omega$, implies that $b \in U$, a contradiction as $S \not\subseteq b$.
\end{proof}

\question{}
\subquestion{}
Let $\Nn = (\NN; R^\Nn)$ be graph such that $(a, b) \in R^\Nn \iff |b - a| = 1$.
Let $D$ be a non-principle ultrafilter on $\NN$ (An example of one is the construction from exercise 2).
Let $\Mm = \Nn^\NN / D$.
We will prove that $\Mm$ is not path-connected.
\begin{proof}
	TODO
\end{proof}

\end{document}
